\chapter{Trésorerie}
\label{ch:tresorerie}

Le module \textbf{Trésorerie} offre une vision globale de votre cash~: factures, soldes bancaires, rapprochement et dépenses.

\section{Synthèse}

\menu{Trésorerie} → \texttt{/treasury}. L'écran de synthèse combine~:

\begin{itemize}[leftmargin=*]
  \item un \textbf{histogramme} mensuel (émis, payé, planifié)~;
  \item un \textbf{tableau} des montants TTC par mois~;
  \item le \textbf{dernier solde bancaire} connu.
\end{itemize}

\screenshot[keepaspectratio]{06-treasury.png}{Synthèse trésorerie — histogramme et tableau mensuel}

La barre \textbf{planifiée} utilise la couleur d'accent \emph{pending} (orange), alignée sur les montants non facturés des écoles. Les montants planifiés du tableau sont exprimés en \textbf{TTC}.

\section{Onglets du module}

\begin{description}[style=nextline, leftmargin=2.5cm]
  \item[Factures] Liste filtrable, statut payé/impayé, accès PDF.
  \item[Créer facture] Émission hors contexte école (cas ponctuels).
  \item[Banque] Banques, comptes, imports de relevés \texttt{.xlsx}.
  \item[Notes de frais] Saisie et validation des frais remboursables.
  \item[Dépenses] Charges autonomes (hors notes de frais).
\end{description}

\section{Import bancaire}

\begin{workflowstep}[Configurer la banque]
  Créez une banque, puis un compte rattaché. Définissez éventuellement le \textbf{compte bancaire de facturation} de l'entreprise pour les KPIs de solde.
\end{workflowstep}

\begin{workflowstep}[Importer un relevé]
  Téléversez un fichier Excel (\texttt{.xlsx}, max 10\,Mo) avec colonnes Date, libellé, débit euros, crédit euros. Vérifiez les lignes sur l'écran de détail d'import.
\end{workflowstep}

\section{Rapprochement}

Pour chaque ligne non lettrée~:

\begin{enumerate}[label=\textbf{\arabic*.}]
  \item Consultez les \textbf{suggestions automatiques} (montant à 0{,}02\,€ près, fenêtre de 45 jours).
  \item Rapprochez avec une \textbf{facture}, un \textbf{groupe de factures} (même école), une \textbf{dépense} ou une \textbf{note de frais}.
  \item Validez~: la facture reçoit \texttt{paid\_at} à la date bancaire si vide.
\end{enumerate}

\begin{avertissement}
  Le \textbf{dérapprochement} supprime le lien mais ne remet pas automatiquement à zéro les dates de paiement~; corrigez manuellement si nécessaire.
\end{avertissement}

\section{Checklist opérationnelle}

\begin{itemize}[leftmargin=*]
  \item Renseigner le solde d'ouverture du compte avant de s'appuyer sur les KPIs.
  \item Utiliser le rapprochement multi-factures uniquement pour un paiement groupé d'une même école.
  \item Vérifier montant, date et libellé avant de valider une suggestion large.
\end{itemize}

\section{Parcours complet rappel}

\begin{center}
\small
\begin{tabular}{@{}cl@{}}
  \toprule
  & \textbf{Action} \\
  \midrule
  1 & Planifier les sessions (Agenda) \\
  2 & Préparer la facturation (fiche école) \\
  3 & Émettre le PDF facture \\
  4 & Importer le relevé bancaire \\
  5 & Rapprocher le crédit → facture payée \\
  \bottomrule
\end{tabular}
\end{center}
